% Lösungen Einheit 2 von 4 – Normalenvektor und Koordinatengleichung (zusammenbau v0.1)
\einheitenkopf{Einheit 2 von 4 · Normalenvektor und Koordinatengleichung}

\erg{12}{a) Koordinatenform \quad b) $\vec n = (-3 | -2 | 1)$ (jedes Vielfache ebenso) \quad c) $d = 2 \cdot 3 + 3 \cdot 2 = 12$ ($y$ kommt nicht vor); Kontrolle mit $C$: $0 + 12 = 12$ ✓ \quad d) $\overrightarrow{CD} = (-3 | 3 | 0)$, $\overrightarrow{CS} = (-3 | -3 | 5)$, $\vec n = (5 | 5 | 6)$; $d = 5 + 5 + 30 = 40$ aus $S$, Probe mit $C$: $20 + 20 = 40$ ✓ \quad e) $\vec n = (6 | 3 | 4)$ aus den Spannvektoren; $d = 24 + 3 = 27$: $E\colon 6x + 3y + 4z = 27$ \quad f) aus $F$: $c = 6$, aus $S$: $a = 3$, aus $M$: $9 + 3b + 12 = 18$, $b = -1$; $W\colon 3x - y + 6z = 18$ \quad g) $\vec x = (9 | 1 | 3) + r \cdot (-8 | 0 | 2) + s \cdot (-6 | 8 | 2)$, $r, s \in \mathbb{R}$; $\vec n = (4 | -1 | 16)$, $d = 36 - 1 + 48 = 83$ aus $W$ (einem Punkt der Plane): $E\colon 4x - y + 16z = 83$ \quad h) $\vec n = (0 | 1 | 0)$, weil beide Richtungsvektoren die $y$-Koordinate 0 haben; mit dem Schnittpunkt: $E\colon y = 3$ (nicht $y = 0$) \quad i) Ansatz $ax + by + cz = 9$: aus $B$ $a = 1$, aus $C$ $b = 3$, aus $D_k$ $c = \frac{9}{k}$: $L_k\colon x + 3y + \frac{9}{k} z = 9$}
\erg{13}{a) Parameterform \quad b) $(4 | -2 | 1)$ und z. B. $(-8 | 4 | -2)$ \quad c) $-3x + y + 4z = -6 + 5 + 4$, also $F\colon -3x + y + 4z = 3$ \quad d) $\overrightarrow{EF} = (-4 | 2 | 0)$, $\overrightarrow{ES} = (-4 | -1 | 2)$; $\vec n = (1 | 2 | 3)$; Probe: $-4 + 4 + 0 = 0$ und $-4 - 2 + 6 = 0$ \quad e) Stützpunkt: $2 + 6 = 8$ ✓, $P$: $4 + 4 = 8$ ✓, Richtungsvektor: $(-4 | 2 | -1) \cdot (0 | 1 | 2) = 2 - 2 = 0$, also liegt ganz $g$ in $H$ \quad f) $(1 | -1 | 0)$ ist ein Vielfaches von $\overrightarrow{BC} = (-3 | 3 | 0)$, $(-4 | 2 | 6)$ das Doppelte von $\overrightarrow{BS} = (-2 | 1 | 3)$; $\vec n$ steht also senkrecht auf zwei nicht parallelen Vektoren der Ebene und damit auf $E$.}
\erg{14}{a) Der Ursprung liegt nicht in $E$; die Konstante kommt aus einem Ebenenpunkt: mit $A$ ist $d = 2 + 2 + 9 = 13$, also $E\colon x + 2y + 3z = 13$ \quad b) Das Skalarprodukt von $k \cdot \vec n$ mit jedem Spannvektor ist das $k$-Fache von null, also null; $k \cdot \vec n$ steht ebenso senkrecht auf der Ebene. \quad c) $E\colon -3y + 5z = 15$; mit $y = 2$: $5z = 21$, $z = 4{,}2$ – das Dach liegt dort 4,2 m hoch \quad d) z. B. $\left(\vec x - (0 | 0 | 5)\right) \cdot (2 | -3 | 1) = 0$}
